\documentclass[a4paper,12pt]{article}

% -------------------------------------------------
% Кодировка, язык, математика
% -------------------------------------------------
\usepackage[T2A]{fontenc}
\usepackage[utf8]{inputenc}
\usepackage[russian]{babel}
\usepackage{amsmath,amssymb,amsthm,mathtools}
\usepackage{icomma}

% -------------------------------------------------
% Типографика
% -------------------------------------------------
\usepackage{helvet}
\renewcommand{\familydefault}{\sfdefault}
\usepackage{microtype}
\usepackage{setspace}
\onehalfspacing
\usepackage{parskip}
\usepackage{enumitem}

% -------------------------------------------------
% Верстка
% -------------------------------------------------
\usepackage[
    a4paper,
    left=22mm,
    right=22mm,
    top=22mm,
    bottom=24mm
]{geometry}

\usepackage{tabularx,booktabs,longtable}
\usepackage{xcolor}

% -------------------------------------------------
% Графика
% -------------------------------------------------
\usepackage{tikz}
\usepackage{tkz-euclide}
\usepackage{pgfplots}
\pgfplotsset{compat=1.18}

% -------------------------------------------------
% Ссылки и колонтитулы
% -------------------------------------------------
\usepackage{hyperref}
\usepackage{fancyhdr}

% -------------------------------------------------
% Палитра
% -------------------------------------------------
\definecolor{accent}{HTML}{008080}
\definecolor{darkaccent}{HTML}{004D4D}
\definecolor{textgray}{HTML}{333333}
\definecolor{softgray}{HTML}{707070}
\definecolor{lightaccent}{HTML}{E8F2F2}

\hypersetup{
    colorlinks=true,
    linkcolor=darkaccent,
    urlcolor=accent
}

\color{textgray}

% -------------------------------------------------
% Служебные настройки
% -------------------------------------------------
\setlength{\parindent}{0pt}
\setlength{\headheight}{15pt}
\renewcommand{\arraystretch}{1.3}

\setlist[itemize]{
    leftmargin=1.5em,
    itemsep=0.25em,
    topsep=0.35em
}

\setlist[enumerate]{
    leftmargin=1.7em,
    itemsep=0.45em,
    topsep=0.4em
}

\pagestyle{fancy}
\fancyhf{}
\fancyhead[L]{\small\textcolor{softgray}{ЕГЭ по математике}}
\fancyhead[R]{\small\textcolor{softgray}{Николь Саркисьянц}}
\fancyfoot[C]{\small\thepage}
\renewcommand{\headrulewidth}{0.3pt}
\renewcommand{\headrule}{
    \hbox to\headwidth{
        \color{accent!55}\leaders\hrule height \headrulewidth\hfill
    }
}

\newcommand{\checktitle}[1]{%
    \vspace{0.4em}
    {\Large\bfseries\textcolor{darkaccent}{#1}}\par
    \vspace{0.25em}
    {\color{accent!55}\rule{\linewidth}{0.6pt}}
    \vspace{0.45em}
}

\newcommand{\subcheck}[1]{%
    \vspace{0.25em}
    {\large\bfseries\textcolor{accent}{#1}}\par
    \vspace{0.15em}
}

\newcommand{\important}[1]{\textbf{\textcolor{darkaccent}{#1}}}

\DeclareMathOperator{\tg}{tg}
\DeclareMathOperator{\ctg}{ctg}
\DeclareMathOperator{\arctg}{arctg}

\begin{document}

% =================================================
% ТИТУЛЬНЫЙ ЛИСТ
% =================================================
\begin{titlepage}
\thispagestyle{empty}

\begin{center}

\vspace*{2.3cm}

{\Large\textcolor{accent}{\textbf{ЧЕК-ЛИСТ}}}

\vspace{0.7cm}

{\huge\bfseries
Тригонометрические уравнения\\[0.25cm]
и смешанные неравенства
}

\vspace{0.6cm}

{\Large\textcolor{darkaccent}{
Преобразования, факторизация,\\
замена переменной и метод интервалов
}}

\vspace{1.3cm}

{\large
Повторение задач повышенного уровня ЕГЭ
}

\vfill

{\large\bfseries
Лёвин Артём Александрович\\[0.15cm]
эксклюзивно для Николь Саркисьянц
}

\vspace{0.7cm}

{\large \today}

\vfill

\begin{tikzpicture}
    \draw[
        accent!55,
        line width=1.1pt
    ]
    (0,0)
    .. controls (2.8,0.35) and (5.5,-0.35) .. (8.1,0)
    .. controls (10.7,0.35) and (13.4,-0.25) .. (16.0,0.10);
\end{tikzpicture}

\vspace{0.4cm}

\end{center}
\end{titlepage}

% =================================================
% ОСНОВНОЙ ЧЕК-ЛИСТ
% =================================================

\checktitle{1. Главная стратегия в тригонометрическом уравнении}

При встрече с нетипичным тригонометрическим уравнением удобно проверять преобразования в следующем порядке:

\begin{enumerate}
    \item \important{Нормализуйте уравнение}: перенесите все слагаемые в одну часть и получите
    \[
        F(x)=0.
    \]

    \item \important{Унифицируйте углы}. Если встречаются \(x\), \(2x\), \(\pi-2x\) и т.\,п., сначала постарайтесь привести выражения к одному аргументу.

    \item \important{Унифицируйте функции}. Например, если основная часть выражения записана через \(\sin x\), удобно и остальные функции выразить через \(\sin x\).

    \item После этого ищите:
    \begin{itemize}
        \item общий множитель;
        \item группировку;
        \item формулы сокращённого умножения;
        \item замену \(t=\sin x\), \(t=\cos x\) и т.\,п.
    \end{itemize}

    \item После замены обязательно учитывайте ограничения:
    \[
        t=\sin x \quad\Longrightarrow\quad -1\le t\le1,
    \]
    \[
        t=\cos x \quad\Longrightarrow\quad -1\le t\le1.
    \]

    \item После решения алгебраического уравнения выполните обратную замену и решите полученные простейшие тригонометрические уравнения.
\end{enumerate}

\important{Ключевой принцип:}
\[
\boxed{
\text{одинаковые углы}
\;\longrightarrow\;
\text{одинаковые функции}
\;\longrightarrow\;
\text{алгебраизация}
}
\]

% -------------------------------------------------

\checktitle{2. Основное тригонометрическое тождество}

Главная формула:
\[
    \sin^2 x+\cos^2 x=1.
\]

Отсюда:
\[
    \cos^2x=1-\sin^2x,
    \qquad
    \sin^2x=1-\cos^2x.
\]

\subcheck{Когда замена особенно удобна}

Если в выражении есть \(\sin x\) и \(\cos^2x\), то
\[
    \cos^2x=1-\sin^2x
\]
сразу приводит всё к одной функции.

Аналогично при наличии \(\cos x\) и \(\sin^2x\).

\subcheck{Когда преобразование обычно невыгодно}

Из равенства
\[
    \sin^2x=1-\cos^2x
\]
нельзя получить
\[
    \sin x=\sqrt{1-\cos^2x}
\]
без потери второго знака.

Правильно:
\[
    \sin x=\pm\sqrt{1-\cos^2x}.
\]

Поэтому переход от \(\sin x\) к \(\cos x\) через извлечение квадратного корня часто \important{усложняет}, а не упрощает задачу.

% -------------------------------------------------

\checktitle{3. Формулы, которые позволяют унифицировать углы}

\subcheck{Двойной угол}

\[
    \sin2x=2\sin x\cos x.
\]

Для косинуса двойного угла существуют три равноценных формы:
\[
    \cos2x=\cos^2x-\sin^2x,
\]
\[
    \cos2x=2\cos^2x-1,
\]
\[
    \cos2x=1-2\sin^2x.
\]

Выбирайте формулу не механически, а в зависимости от того, \important{к какой функции нужно привести уравнение}.

Например, если остальные слагаемые содержат \(\sin x\), наиболее естественна формула
\[
    \cos2x=1-2\sin^2x.
\]

\subcheck{Сумма и разность углов}

\[
    \cos(\alpha-\beta)
    =
    \cos\alpha\cos\beta
    +
    \sin\alpha\sin\beta,
\]

\[
    \cos(\alpha+\beta)
    =
    \cos\alpha\cos\beta
    -
    \sin\alpha\sin\beta,
\]

\[
    \sin(\alpha-\beta)
    =
    \sin\alpha\cos\beta
    -
    \cos\alpha\sin\beta,
\]

\[
    \sin(\alpha+\beta)
    =
    \sin\alpha\cos\beta
    +
    \cos\alpha\sin\beta.
\]

Полезный частный случай:
\[
\begin{aligned}
    \cos(\pi-2x)
    &=
    \cos\pi\cos2x+\sin\pi\sin2x\\
    &=
    -\cos2x.
\end{aligned}
\]

\important{На экзамене формулы могут быть в справочном материале, но знание их наизусть помогает прежде всего распознать момент применения.}

% -------------------------------------------------

\checktitle{4. Группировка и двойное вынесение}

Если после преобразований получилось несколько слагаемых, ищите разбиение на группы так, чтобы после первого вынесения появились \important{одинаковые скобки}.

Пример:
\[
    2t^3+\sqrt2\,t^2-2t-\sqrt2=0.
\]

Группируем:
\[
    \left(2t^3+\sqrt2\,t^2\right)
    +
    \left(-2t-\sqrt2\right)=0.
\]

Выносим множители:
\[
    t^2(2t+\sqrt2)-(2t+\sqrt2)=0.
\]

Повторно выносим общую скобку:
\[
    (2t+\sqrt2)(t^2-1)=0.
\]

И окончательно:
\[
    (2t+\sqrt2)(t-1)(t+1)=0.
\]

Это и есть схема \important{двойного вынесения}:
\[
\boxed{
\text{группировка}
\to
\text{первое вынесение}
\to
\text{одинаковые скобки}
\to
\text{второе вынесение}
}
\]

\subcheck{Полезная эвристика}

Чётное число слагаемых действительно часто подсказывает возможность группировки, однако это \important{не теорема}.

Правильная формулировка:

\begin{itemize}
    \item четыре или шесть слагаемых --- повод проверить группировку;
    \item нечётное число слагаемых --- повод искать сокращение, замену или формулу;
    \item окончательное решение определяется структурой выражения, а не числом его слагаемых.
\end{itemize}

% -------------------------------------------------

\checktitle{5. Пример: полная схема решения тригонометрического уравнения}

Решим
\[
    2\sin^3x
    =
    \sqrt2\,\cos^2x+2\sin x.
\]

Переносим всё в левую часть:
\[
    2\sin^3x-\sqrt2\,\cos^2x-2\sin x=0.
\]

Так как присутствуют \(\sin x\) и \(\cos^2x\), заменяем
\[
    \cos^2x=1-\sin^2x.
\]

Важно сохранить скобки:
\[
    2\sin^3x-\sqrt2(1-\sin^2x)-2\sin x=0.
\]

После раскрытия:
\[
    2\sin^3x+\sqrt2\,\sin^2x-2\sin x-\sqrt2=0.
\]

Пусть
\[
    t=\sin x,
    \qquad
    -1\le t\le1.
\]

Тогда
\[
    2t^3+\sqrt2\,t^2-2t-\sqrt2=0.
\]

Группировка:
\[
    t^2(2t+\sqrt2)-(2t+\sqrt2)=0,
\]
\[
    (2t+\sqrt2)(t^2-1)=0.
\]

Отсюда
\[
    t=-\frac{\sqrt2}{2},
    \qquad
    t=1,
    \qquad
    t=-1.
\]

Все значения удовлетворяют условию
\[
    -1\le t\le1.
\]

Получаем:
\[
    \sin x=-\frac{\sqrt2}{2},
    \qquad
    \sin x=1,
    \qquad
    \sin x=-1.
\]

Следовательно,
\[
    x=-\frac{\pi}{4}+2\pi n,
\]
или
\[
    x=\frac{5\pi}{4}+2\pi n,
\]
или
\[
    x=\frac{\pi}{2}+2\pi n,
\]
или
\[
    x=-\frac{\pi}{2}+2\pi n,
    \qquad n\in\mathbb Z.
\]

% -------------------------------------------------

\checktitle{6. Простейшие тригонометрические уравнения}

\subcheck{Синус}

Для
\[
    \sin x=a,
    \qquad |a|\le1,
\]
удобно использовать две серии:
\[
    x=\arcsin a+2\pi n,
\]
\[
    x=\pi-\arcsin a+2\pi n,
    \qquad n\in\mathbb Z.
\]

Эквивалентная компактная формула:
\[
    x=(-1)^n\arcsin a+\pi n,
    \qquad n\in\mathbb Z.
\]

\subcheck{Косинус}

Для
\[
    \cos x=a,
    \qquad |a|\le1,
\]
получаем
\[
    x=\pm\arccos a+2\pi n,
    \qquad n\in\mathbb Z.
\]

\subcheck{Частные случаи}

\[
\begin{array}{ccl}
\sin x=1
&\Longrightarrow&
x=\dfrac{\pi}{2}+2\pi n,\\[0.3cm]

\sin x=-1
&\Longrightarrow&
x=-\dfrac{\pi}{2}+2\pi n,\\[0.3cm]

\sin x=0
&\Longrightarrow&
x=\pi n,\\[0.3cm]

\cos x=1
&\Longrightarrow&
x=2\pi n,\\[0.3cm]

\cos x=-1
&\Longrightarrow&
x=\pi+2\pi n,
\end{array}
\qquad n\in\mathbb Z.
\]

% -------------------------------------------------

\checktitle{7. Упрощение квадратных корней}

Перед группировкой полезно максимально упростить радикалы.

Например:
\[
    \sqrt8=2\sqrt2,
    \qquad
    \sqrt{12}=2\sqrt3,
    \qquad
    \sqrt6=\sqrt2\sqrt3.
\]

Это позволяет увидеть повторяющиеся элементы.

Например:
\[
    4t^2+\sqrt8\,t-\sqrt6-\sqrt{12}\,t=0
\]
превращается в
\[
    4t^2+2\sqrt2\,t-\sqrt6-2\sqrt3\,t=0.
\]

После группировки:
\[
    2t(2t+\sqrt2)-\sqrt3(2t+\sqrt2)=0,
\]
откуда
\[
    (2t+\sqrt2)(2t-\sqrt3)=0.
\]

\important{Идея: большие подкоренные числа часто скрывают общие множители.}

% -------------------------------------------------

\checktitle{8. Формулы сокращённого умножения, которые необходимо узнавать}

\subcheck{Квадрат разности}

\[
    (a-b)^2=a^2-2ab+b^2.
\]

\subcheck{Квадрат суммы}

\[
    (a+b)^2=a^2+2ab+b^2.
\]

\subcheck{Куб разности}

\[
    (a-b)^3
    =
    a^3-3a^2b+3ab^2-b^3.
\]

\subcheck{Куб суммы}

\[
    (a+b)^3
    =
    a^3+3a^2b+3ab^2+b^3.
\]

\subcheck{Разность кубов}

\[
    a^3-b^3
    =
    (a-b)(a^2+ab+b^2).
\]

\subcheck{Сумма кубов}

\[
    a^3+b^3
    =
    (a+b)(a^2-ab+b^2).
\]

Особенно полезно узнавать запись
\[
    t^3-9t^2+27t-27
\]
как
\[
    t^3-3t^2\cdot3+3t\cdot3^2-3^3,
\]
то есть
\[
    \boxed{(t-3)^3}.
\]

\subcheck{Важное уточнение о теореме Безу}

Теорема Безу утверждает:
\[
    P(a)=0
    \quad\Longleftrightarrow\quad
    P(x)\ \text{делится на}\ x-a.
\]

Она сама по себе не является правилом поиска корней по целочисленным делителям.

Поиск рациональных корней многочлена с целыми коэффициентами связан с отдельным критерием о возможных рациональных корнях.

Поэтому при иррациональных коэффициентах не следует механически пытаться применять «делители свободного члена».

% -------------------------------------------------

\checktitle{9. Смешанное дробно-рациональное неравенство}

Если неравенство уже имеет вид
\[
    \frac{F(x)}{G(x)}>0,
\]
то оно готово к \important{методу интервалов}.

Алгоритм:

\begin{enumerate}
    \item Решите
    \[
        F(x)=0.
    \]

    \item Решите
    \[
        G(x)=0.
    \]

    \item Корни знаменателя обязательно исключите из ОДЗ:
    \[
        G(x)\ne0.
    \]

    \item Расположите все критические точки на числовой прямой.

    \item Определите знак выражения на каждом интервале.

    \item Выберите интервалы в соответствии со знаком исходного неравенства.
\end{enumerate}

\important{Нельзя делать вывод «знаменатель не дал подходящих корней, значит неравенство не имеет решений».}

Знаменатель прежде всего определяет:
\begin{itemize}
    \item запрещённые точки;
    \item знак дроби.
\end{itemize}

% -------------------------------------------------

\checktitle{10. Показательная часть: сначала создайте повторяющийся элемент}

Рассмотрим выражение
\[
    27^x-9^{x+1}+3^{x+3}-27.
\]

Приведём всё к основанию \(3\):
\[
    27^x=3^{3x}=(3^x)^3,
\]
\[
    9^{x+1}
    =
    9\cdot9^x
    =
    9(3^x)^2,
\]
\[
    3^{x+3}=27\cdot3^x.
\]

Пусть
\[
    t=3^x.
\]

Так как показательная функция положительна,
\[
    \boxed{t>0}.
\]

Получаем:
\[
    t^3-9t^2+27t-27.
\]

Узнаём куб разности:
\[
    t^3-9t^2+27t-27=(t-3)^3.
\]

Поэтому
\[
    t=3
    \quad\Longrightarrow\quad
    3^x=3
    \quad\Longrightarrow\quad
    x=1.
\]

\important{Типовая схема:}
\[
\boxed{
\text{одинаковое основание}
\to
\text{повторяющийся элемент}
\to
\text{замена}
\to
\text{алгебраическое уравнение}
}
\]

% -------------------------------------------------

\checktitle{11. Кратность корня и смена знака}

Если выражение содержит
\[
    (x-a)^k,
\]
то число \(k\) называется кратностью корня \(a\).

При прохождении через корень:

\[
\begin{array}{ccl}
k\ \text{нечётно}
&\Longrightarrow&
\text{знак меняется},\\[0.2cm]

k\ \text{чётно}
&\Longrightarrow&
\text{знак не меняется}.
\end{array}
\]

Например:
\[
    (t-3)^3
\]
имеет корень \(t=3\) нечётной кратности, поэтому знак меняется.

А выражение
\[
    (x-a)^2
\]
сохраняет знак по обе стороны от \(a\).

Если корень принадлежит \important{знаменателю}, точка всё равно исключается из области допустимых значений независимо от кратности.

% -------------------------------------------------

\checktitle{12. Пример полного анализа дробного неравенства}

Рассмотрим
\[
    \frac{
        27^x-9^{x+1}+3^{x+3}-27
    }{
        50x^2-110x+\dfrac{121}{2}
    }>0.
\]

Числитель:
\[
    27^x-9^{x+1}+3^{x+3}-27
    =
    (3^x-3)^3.
\]

Его единственный корень:
\[
    x=1,
\]
причём кратность равна \(3\).

Теперь знаменатель:
\[
    50x^2-110x+\frac{121}{2}.
\]

Умножим уравнение знаменателя на \(2\):
\[
    100x^2-220x+121=0.
\]

Дискриминант:
\[
\begin{aligned}
    D
    &=
    (-220)^2-4\cdot100\cdot121\\
    &=
    48400-48400\\
    &=0.
\end{aligned}
\]

Следовательно,
\[
    x=\frac{220}{200}=\frac{11}{10}.
\]

Более того,
\[
    50x^2-110x+\frac{121}{2}
    =
    50\left(x-\frac{11}{10}\right)^2.
\]

Поэтому знаменатель:
\begin{itemize}
    \item положителен при \(x\ne\dfrac{11}{10}\);
    \item равен нулю при \(x=\dfrac{11}{10}\).
\end{itemize}

Критические точки:
\[
    1,
    \qquad
    \frac{11}{10}.
\]

\begin{center}
\begin{tikzpicture}[x=1cm,y=1cm]
    \draw[->,line width=0.9pt] (-0.3,0) -- (10.2,0);

    \draw[line width=0.8pt] (3,-0.10) -- (3,0.10);
    \filldraw[fill=white,draw=darkaccent,line width=0.9pt]
        (3,0) circle (2.5pt);
    \node[below=6pt] at (3,0) {$1$};

    \draw[line width=0.8pt] (6.4,-0.10) -- (6.4,0.10);
    \filldraw[fill=white,draw=darkaccent,line width=0.9pt]
        (6.4,0) circle (2.5pt);
    \node[below=6pt] at (6.4,0) {$\dfrac{11}{10}$};

    \node[above=5pt] at (1.3,0) {$-$};
    \node[above=5pt] at (4.7,0) {$+$};
    \node[above=5pt] at (8.2,0) {$+$};

    \draw[accent,line width=2pt] (3.15,0.35) -- (6.25,0.35);
    \draw[accent,line width=2pt] (6.55,0.35) -- (9.7,0.35);
\end{tikzpicture}
\end{center}

Ответ:
\[
    \boxed{
    x\in
    \left(1;\frac{11}{10}\right)
    \cup
    \left(\frac{11}{10};+\infty\right)
    }.
\]

% -------------------------------------------------

\checktitle{13. Типичные ошибки}

\begin{enumerate}
    \item \important{Потеря скобок при подстановке.}

    Неверно:
    \[
        -\sqrt2\cos^2x
        \to
        -\sqrt2-\sin^2x.
    \]

    Правильно:
    \[
        -\sqrt2(1-\sin^2x).
    \]

    \item \important{Попытка решить всё как квадратное уравнение}, хотя выражение специально устроено под группировку или формулу сокращённого умножения.

    \item \important{Игнорирование ограничения после замены}
    \[
        t=\sin x
        \quad\Longrightarrow\quad
        -1\le t\le1.
    \]

    \item \important{Неупрощённые радикалы}:
    \[
        \sqrt8,\ \sqrt{12}
    \]
    могут скрывать общие множители.

    \item \important{Неверный вывод из дискриминанта знаменателя.}

    Даже если знаменатель не имеет действительных корней, необходимо установить его знак.

    Если \(D=0\), корень существует и является запрещённой точкой.

    \item \important{Сокращение множителя со знаменателем без фиксации ОДЗ.}

    Даже после алгебраического сокращения исходная запрещённая точка остаётся запрещённой.

    \item \important{Автоматическое чередование знаков} при методе интервалов.

    Знак меняется только при прохождении через корень нечётной кратности.

    \item \important{Ошибка в общей формуле простейшего тригонометрического уравнения} после правильно выполненной основной части решения.

    \item \important{Преждевременное вычисление дискриминанта} громоздкого многочлена вместо поиска структуры и факторизации.
\end{enumerate}

% -------------------------------------------------

\checktitle{14. Финальная памятка перед решением}

Перед тем как переходить к вычислениям, задайте себе вопросы:

\begin{enumerate}
    \item Всё ли перенесено в одну часть?
    \item Одинаковы ли аргументы тригонометрических функций?
    \item Можно ли сделать функции одинаковыми?
    \item Есть ли формула двойного угла?
    \item Есть ли формула суммы или разности углов?
    \item Можно ли сделать замену?
    \item Выписано ли ограничение на новую переменную?
    \item Можно ли сгруппировать слагаемые?
    \item Не скрывается ли формула сокращённого умножения?
    \item Упрощены ли квадратные корни?
    \item Если есть дробь, найдены ли нули числителя и знаменателя отдельно?
    \item Все ли нули знаменателя исключены?
    \item Учтена ли кратность каждого корня?
\end{enumerate}

\[
\boxed{
\text{Сначала структура задачи --- потом вычисления.}
}
\]

% =================================================
% ТРЕНИРОВОЧНЫЙ БЛОК
% =================================================
\newpage

\thispagestyle{fancy}

\begin{center}
    {\LARGE\bfseries\textcolor{darkaccent}{Тренировочный блок}}
    
    \vspace{0.2cm}
    
    {\large 10 упражнений по нарастающей сложности}
\end{center}

\vspace{0.5cm}

\subcheck{Средний уровень}

\begin{enumerate}[label=\textbf{\arabic*.}]

\item Решите уравнение:
\[
    2\sin^3x
    =
    \sqrt2\,\cos^2x+2\sin x.
\]

\item Решите уравнение:
\[
    \sin2x+\sqrt2\,\sin x
    =
    2\cos x+\sqrt2.
\]

\item Решите уравнение:
\[
    \cos(\pi-x)=\frac12.
\]

\end{enumerate}

\subcheck{Выше среднего}

\begin{enumerate}[label=\textbf{\arabic*.},start=4]

\item Решите уравнение:
\[
    2\cos(\pi-2x)+2
    +\sqrt8\,\sin x
    -\sqrt{12}\,\sin x
    -\sqrt6
    =0.
\]

\item Решите уравнение:
\[
    \cos2x+\sin x=0.
\]

\item Решите уравнение:
\[
    \sin2x=\sqrt3\,\sin x.
\]

\item Решите неравенство:
\[
    \frac{
        27^x-9^{x+1}+3^{x+3}-27
    }{
        50x^2-110x+\dfrac{121}{2}
    }>0.
\]

\item Решите неравенство:
\[
    \frac{(2^x-2)^3}{(x-1)(x-3)}
    \le0.
\]

\item Решите уравнение:
\[
    8^x-12\cdot4^x+48\cdot2^x-64=0.
\]

\end{enumerate}

\subcheck{Сложный уровень}

\begin{enumerate}[label=\textbf{\arabic*.},start=10]

\item Решите неравенство:
\[
    \frac{
        27^x-9^{x+1}+3^{x+3}-27
    }{
        9^x-10\cdot3^x+9
    }>0.
\]

\end{enumerate}

% =================================================
% ОТВЕТЫ
% =================================================
\newpage

\begin{center}
    {\LARGE\bfseries\textcolor{darkaccent}{Ответы}}
\end{center}

\vspace{0.5cm}

\begin{table}[ht]
\centering
\caption{Ответы к тренировочному блоку}
\renewcommand{\arraystretch}{1.45}
\begin{tabularx}{\linewidth}{@{}>{\bfseries}p{1.1cm}X@{}}
\toprule
№ & Ответ \\
\midrule

1 &
\(
x=-\dfrac{\pi}{4}+2\pi n;
\quad
x=\dfrac{5\pi}{4}+2\pi n;
\quad
x=\dfrac{\pi}{2}+2\pi n;
\quad
x=-\dfrac{\pi}{2}+2\pi n,
\quad n\in\mathbb Z.
\)
\\

2 &
\(
x=\dfrac{\pi}{2}+2\pi n;
\quad
x=\dfrac{3\pi}{4}+2\pi n;
\quad
x=\dfrac{5\pi}{4}+2\pi n,
\quad n\in\mathbb Z.
\)
\\

3 &
\(
x=\dfrac{2\pi}{3}+2\pi n
\)
или
\(
x=\dfrac{4\pi}{3}+2\pi n,
\quad n\in\mathbb Z.
\)
\\

4 &
\(
x=\dfrac{\pi}{3}+2\pi n;
\quad
x=\dfrac{2\pi}{3}+2\pi n;
\quad
x=-\dfrac{\pi}{4}+2\pi n;
\quad
x=\dfrac{5\pi}{4}+2\pi n,
\quad n\in\mathbb Z.
\)
\\

5 &
\(
x=\dfrac{\pi}{2}+2\pi n;
\quad
x=-\dfrac{\pi}{6}+2\pi n;
\quad
x=\dfrac{7\pi}{6}+2\pi n,
\quad n\in\mathbb Z.
\)
\\

6 &
\(
x=\pi n
\)
или
\(
x=\pm\dfrac{\pi}{6}+2\pi n,
\quad n\in\mathbb Z.
\)
\\

7 &
\(
\displaystyle
x\in
\left(1;\frac{11}{10}\right)
\cup
\left(\frac{11}{10};+\infty\right).
\)
\\

8 &
\(
x\in(-\infty;1)\cup(1;3).
\)
\\

9 &
\(
x=2.
\)
\\

10 &
\(
x\in(0;1)\cup(2;+\infty).
\)
\\

\bottomrule
\end{tabularx}
\end{table}

\vfill

\begin{center}
\textcolor{softgray}{
\small
При проверке решений обращайте особое внимание на ОДЗ,
кратность корней и запрещённые значения знаменателя.
}
\end{center}

\end{document}